\section{Q-expressions and \QED{}}
\label{sec:qexpr}

\emph{Source: Wang, Pan, Cheung, PVLDB 17(11):3602--3614, 2024~\cite{qed}.}

\subsection{The one-line summary}

Q-expressions have \emph{the same syntax} as U-expressions. What changes is that they are
pinned to a single concrete semantic domain, which makes every operator \emph{definable}
rather than axiomatised, which in turn makes them cheaply encodable into SMT. Everything
else in \QED{} follows from that.

\begin{table}[h]\centering\small
\begin{tabularx}{\linewidth}{@{}lXX@{}}
\toprule
 & \textbf{U-expression} & \textbf{Q-expression}\\
\midrule
semantics & any U-semiring; operators fixed by \emph{axioms} & the single model $\Nbar$; operators fixed by \emph{definitions}\\
$\sq{a}$ & 6 axiom lines (7 identities) & \emph{defined} as $\ind{a\neq0}$\\
$\notu{a}$ & 4 axiom lines & \emph{defined} as $\ind{a=0}$\\
decision & syntactic isomorphism & SMT, used at every stage\\
NULL & absent & first class (option types, 3-valued logic)\\
unknown operators & unsupported & $\qop$/$\hop$, first class\\
\bottomrule
\end{tabularx}
\end{table}

\subsection{The concrete domain}

\begin{definition}[\cite{qed} Eq.\ 4]\label{def:nbar}
Multiplicities live in $\Nbar=\mathbb{N}\cup\{\infty\}$ with
\[
\infty+a=\infty,\qquad 0\times\infty=0,\qquad a\times\infty=\infty\ (a\neq0),
\]
\[
\usum{s} f(s)=\begin{cases} f(s_1)+\cdots+f(s_n) & \text{if } \mathrm{supp}(f)=\{s_1,\dots,s_n\}\text{ is finite}\\ \infty & \text{otherwise,}\end{cases}
\]
\[
\ind{P}=\begin{cases}1&P\text{ true}\\0&\text{else,}\end{cases}
\qquad \sq{a}=\ind{a\neq0},\qquad \notu{a}=\ind{a=0}.
\]
A \emph{Q-expression} is any expression of type $\Nbar$.
\end{definition}

\keybox{Why this is more than bookkeeping}{%
$\Nbar$ \emph{is} a U-semiring, so nothing is lost. But because $\sq{a}$ is now literally
$\ind{a\neq0}$, a squash of a table application is just a Boolean --- an uninterpreted
predicate an SMT solver can reason about directly. In the axiomatic setting $\sq{\cdot}$ is
an opaque symbol you may only rewrite with the six squash axioms. This is the difference
between ``call Z3'' and ``search for an axiom chain''.}

\subsection{Query syntax}

\QED{} fixes a small relational core~\cite[Fig.\ 2]{qed}:
\[
\begin{array}{r@{\;}l@{\qquad}r@{\;}l}
Q ::=& \texttt{Table}(R:\sigma) & \mid& \texttt{Values}(v_1,\dots,v_n)\\
\mid& \texttt{Filter}(P,Q) & \mid& \texttt{Proj}(f,Q)\\
\mid& \texttt{Join}(Q_1,Q_2) & \mid& \texttt{Union}(Q_1,Q_2)\\
\mid& \texttt{Minus}(Q_1,Q_2) & \mid& \texttt{Distinct}(Q)\\
\mid& \texttt{GroupBy}(k,\alpha(f),Q) & \mid& \qop(o,v_1,\dots,v_n,Q)\\[4pt]
v,P,f ::=& x \mid \texttt{Exists}(Q) & \mid& \texttt{Op}(o,\vec v) \mid \hop(o,\vec v,Q)
\end{array}
\]
Two entries deserve comment.

$\qop$ and $\hop$ are \textbf{uninterpreted operators} --- \emph{first-class} escape
hatches. \lstinline|LIMIT 10| is $\qop(\texttt{Limit},10,Q)$; \lstinline|a IN Q| is
$\hop(\texttt{In},a,Q)$. Anything \QED{} does not model becomes one of these and still
participates in reasoning through congruence, rather than aborting the proof.

\begin{remark}[what is deliberately missing]
Set \texttt{Minus} and set \texttt{Intersect} are modelled; \emph{bag} minus and
\emph{bag} intersect are not. The paper's footnote is explicit: those would need a
\emph{monus} (truncated subtraction) at the semiring level, which \QED{} declines to add.
Handle them as $\qop$.
\end{remark}

\subsection{Semantics}
\label{sec:qsem}

\begin{align*}
\den{\texttt{Table}(R{:}\sigma)} &= \lam{x}{R(x)}\\
\den{\texttt{Values}(v_1..v_n)} &= \lam{x}{\ind{x{=}v_1}+\cdots+\ind{x{=}v_n}}\\
\den{\texttt{Filter}(P,Q)} &= \lam{x}{\ind{\den{P}(x)=\textsf{True}}\times\den{Q}(x)}\\
\den{\texttt{Proj}(f,Q)} &= \lam{x}{\usum{s}\ind{x=\den{f}(s)}\times\den{Q}(s)}\\
\den{\texttt{Join}(Q_1,Q_2)} &= \lam{x_1,x_2}{\den{Q_1}(x_1)\times\den{Q_2}(x_2)}\\
\den{\texttt{Union}(Q_1,Q_2)} &= \lam{x}{\den{Q_1}(x)+\den{Q_2}(x)}\\
\den{\texttt{Distinct}(Q)} &= \lam{x}{\sq{\den{Q}(x)}}\\
\den{\texttt{Minus}(Q_1,Q_2)} &= \lam{x}{\sq{Q_1(x)\times\notu{Q_2(x)}}}\\
\den{\texttt{GroupBy}(k,\alpha(f),Q)} &= \lam{x,y}{
   \underbrace{\Bigl\lVert\usum{s}\ind{x=\den{k}(s)}\times\den{Q}(s)\Bigr\rVert}_{\text{the group is non-empty}}
   \times \underbrace{\Bigl[\,y=\hop\bigl(\alpha,\lam{y'}{\usum{s}\den{Q}(s)\times\ind{x=\den{k}(s)\wedge y'=\den{f}(s)}}\bigr)\Bigr]}_{\text{$y$ is the aggregate of the group's bag}}}
\end{align*}
and for scalars, $\den{\texttt{Exists}(Q)}=\sq{\usum{s}\den{Q}(s)}$.

\paragraph{Read the \texttt{GroupBy} rule slowly --- it is the crux of the paper.}
The output row is the pair $(x,y)$ = (group key, aggregate value). Two factors:
\begin{itemize}
\item The \emph{squashed} sum says ``at least one input row has key $x$'', and the squash is
      what makes \lstinline|GROUP BY| emit \textbf{one} row per key regardless of input
      multiplicities. (This is exactly the subtlety \udp{}'s de-sugaring glosses over ---
      Remark~\ref{rem:groupby}.)
\item The bracket pins $y$ to $\alpha$ applied to a \emph{bag}: the function
      $\lam{y'}{\dots}$ whose value at $y'$ is the multiplicity of $y'$ within the group.
      $\alpha$ is an $\hop$, i.e.\ \textbf{uninterpreted}. \QED{} knows only that equal bags
      give equal aggregates.
\end{itemize}

\begin{definition}[the equivalence problem]
Write $\Gamma\mid\Delta\mid\Phi \vdash Q_1 \stackrel{?}{=} Q_2$, where $\Gamma$ are table
symbols, $\Delta$ other uninterpreted symbols and $\Phi$ an assumed predicate. It asks
whether $\den{Q_1}(s)=\den{Q_2}(s)$ for all instantiations of $\Gamma,\Delta$ satisfying
$\Phi$ and all rows $s$.
\end{definition}

\subsection{Integrity constraints: rewrite the table, not the query}

\paragraph{Primary key (\cite{qed} Eqs.\ 6--8).}
Let $R$ have schema $K\times S$ with key column(s) $k$. A table with a primary key is
\emph{exactly} a subset of $K$ plus a function $f_R:K\to S$ giving the non-key columns.
Formally, for any predicate $P$ over tables,
\begin{equation}\label{eq:pk-complete}
\forall R \text{ with primary key } k.\ P(R)
\iff \forall R',f_R.\ P\bigl(\lam{k,s}{\sq{R'(k)}\times\ind{s=f_R(k)}}\bigr)
\end{equation}
which licenses the rewrite
\begin{equation}\label{eq:pk7}
R \;\rightsquigarrow\; \lam{k,s}{\sq{R'(k)}\times\ind{s=f_R(k)}}
\end{equation}
introducing a fresh table symbol $R'$ of schema $K$ and a fresh function symbol $f_R$.
When $k$ is the \emph{only} column this collapses to the form everyone quotes:
\begin{equation}\label{eq:pk8}
R \;\rightsquigarrow\; \lam{k}{\sq{R(k)}}.
\end{equation}

\warnbox{Do not quote Eq.~\eqref{eq:pk8} when the table has non-key columns}{%
In our running example $S$ has schema $(a,b)$ with key $a$, so the applicable rule is
\eqref{eq:pk7}: $S\rightsquigarrow\lam{a,b}{\sq{S'(a)}\times\ind{b=f_S(a)}}$.
The familiar $\lam{a}{\sq{S(a)}}$ appears only after $b$ is projected away --- which the
queries do, since neither mentions $S.b$. The distinction matters the moment a query reads
a non-key column.}

The word \emph{complete} in the paper is Eq.~\eqref{eq:pk-complete}: the rewrite yields an
equivalence problem that is equivalent to the original, not merely implied by it. Contrast
\spes{}, whose primary-key support is a finite list of query-level rewrite rules --- sound,
but only firing on shapes someone anticipated.

\paragraph{Foreign key (\cite{qed} Eq.\ 9).}
\[
S \;\rightsquigarrow\; S,R' \mid \forall b,k.\ \sq{S(b,k)}\to\sq{R'(k)} \;\vdash\; S
\]
Note what changed relative to \udp{}: \udp{} applies its foreign-key identity
\emph{eagerly} as a rewrite, which does not terminate when two tables reference each other.
\QED{} emits \emph{one global logical formula} instead and lets the SMT solver unfold its
consequences lazily. That is a strictly better engineering decision and removes \udp{}'s
non-termination hazard.

\paragraph{Check constraint.} $\forall s.\ \sq{R(s)}\to C(s)$. One line, same idea.

\subsection{NULL semantics}

\udp{} has none; this alone rules out outer joins. \QED{} lifts every type $T$ to
$\mathrm{Option}(T)$, an algebraic datatype modern SMT solvers support natively.

\begin{itemize}
\item \textbf{Strict by default:} an $n$-ary operation returns \textsf{Null} exactly when
      some argument is \textsf{Null}. This lifts $+$, $\times$, comparisons, etc.
\item \textbf{Three-valued \texttt{AND}/\texttt{OR}} are the documented exceptions:
\begin{align*}
a\ \texttt{And}\ b &\coloneqq \mathrm{ite}(a=\textsf{Null},\ \mathrm{ite}(b=\textsf{True},\textsf{Null},b),\ \mathrm{ite}(a=\textsf{True},b,a))\\
a\ \texttt{Or}\ b &\coloneqq \mathrm{ite}(a=\textsf{Null},\ \mathrm{ite}(b=\textsf{False},\textsf{Null},b),\ \mathrm{ite}(a=\textsf{False},b,a))
\end{align*}
      Check them: \textsf{Null} \texttt{And} \textsf{False} $=$ \textsf{False} (not
      \textsf{Null}), and \textsf{Null} \texttt{Or} \textsf{True} $=$ \textsf{True}.
      Strictness would get both wrong.
\item \textbf{Quantified comparison:}
\[
a\ \langle\mathrm{cmp}\rangle\ \texttt{Some}\ Q \coloneqq
\mathrm{ite}\bigl(\exists x.\sq{Q(x)}\wedge (x\,\langle\mathrm{cmp}\rangle\,a)=\textsf{True},\ \textsf{True},\
\mathrm{ite}(\exists x.\sq{Q(x)}\wedge(x\,\langle\mathrm{cmp}\rangle\,a)=\textsf{Null},\ \textsf{Null},\ \textsf{False})\bigr)
\]
      \lstinline|a IN Q| is the special case with $\langle\mathrm{cmp}\rangle$ being
      \texttt{=}. The nesting order encodes SQL's rule: \emph{a true match wins over a
      null one}, and only if there is no match at all and no null do you get
      \textsf{False}.
\item \textbf{\lstinline|COUNT(*)| vs \lstinline|COUNT(k)|:} modelled by optionally applying
      a non-\textsf{Null} filter before the aggregate.
\end{itemize}

\subsection{Algorithm A: the complete fragment}
\label{sec:qed-complete}

\QED{} has \emph{two} algorithms and picks between them. The first is complete on a
restricted fragment; this is the paper's theoretical contribution.

\begin{definition}[$\mathcal F_T$, \cite{qed} Def.\ 1]
Fix a first-order theory $T$. The fragment $\mathcal F_T$ consists of
\texttt{Table}, \texttt{Values}, \texttt{Filter}, \texttt{Proj}, \texttt{Join},
\texttt{Union}, where all values, projection functions and filter predicates are definable
in $T$. Primary keys on the tables are allowed.
\end{definition}

This generalises Cohen's classical UCQ result~\cite{cohen}: UCQ forces predicates to be
conjunctions of (in)equalities and projections to be bare column references, whereas
$\mathcal F_T$ allows \emph{anything expressible in $T$}. With $T$ = the SMT solver's
theories, that is a large class. Note $\mathcal F_T$ excludes \texttt{GroupBy},
\texttt{Distinct} and \texttt{Minus} --- so our running example is \emph{not} in it.

\subsubsection{Scoped normal form (SNF)}
\begin{definition}[\cite{qed} Def.\ 2]
$U$ is in SNF if $U=T_1+\cdots+T_n$ where each \emph{normal term} is
\[
T=\usum{R_1^{k_1}(s_1)}\cdots\usum{R_m^{k_m}(s_m)} \ind{P},
\qquad
\usum{R^0(s)}U \coloneqq \usum{s}\sq{R(s)}\times U,\quad
\usum{R^k(s)}U \coloneqq \usum{s}\underbrace{R(s)\cdots R(s)}_{k}\times U
\]
with $P$ (the \emph{body}, $\mathrm{Bdy}\,T$) containing no table variables, and the
summations ordered consistently with a fixed order on table symbols.
\end{definition}
Auxiliary notation used below: $\mathrm{Scp}\,T$ is the list of all summation variables;
$\mathrm{Scp}_R T$ those attached to table $R$; $\mathrm{Rel}\,T$ the list of table
variables with their powers; $\mathrm{Rel}'T$ the same without powers.

\noindent Normalisation (\cite{qed} Alg.\ 1) is a straightforward induction:
\texttt{Table} is already a term; \texttt{Filter} conjoins its predicate into every body;
\texttt{Proj} post-composes $f$ into the body \emph{without} adding a summation;
\texttt{Join} merges scopes and conjoins bodies pairwise; \texttt{Union} concatenates terms.

\subsubsection{Linearisation (LNF)}
Two obstacles remain. First, summation variables attached to the same table can be
introduced in any order, so terms that differ only by a permutation look different.
Second, redundant variables: in
\[
T_1=\usum{R(s_1)}\usum{R(s_2)}\ind{(x{=}s_2)\wedge(s_1{=}s_2)}
\qquad\text{vs}\qquad
T_2=\usum{R^2(s_1)}\ind{x{=}s_1}
\]
$T_1$'s $s_2$ is redundant, but the scopes have different lengths so they cannot be
compared directly.

\begin{definition}[LNF, \cite{qed} Def.\ 3]
$U$ is in LNF if it is in SNF and, for every term $T$ and table variable $R$, $\mathrm{Bdy}\,T$
implies the variables of $\mathrm{Scp}_R T$ are pairwise distinct and form a linearly
ordered chain.
\end{definition}

Algorithm 2 achieves this by \emph{case-splitting on the order}. For a scope
$\{a,b,c\}$, $\mathrm{Lin}$ enumerates all $13$ ways three values can be ordered-or-equal
($a{<}b{<}c$, \dots, $a{=}b{<}c$, \dots, $a{=}b{=}c$). Each term is replaced by the sum over
all cases; within each case, variables proved equal are merged (the oracle decides
$\mathrm{Bdy}\,T'\to s_j=s_k$), and the survivors are sorted ascending. Now every term has a
canonical variable order.

\begin{remark}[cost]
$|\mathrm{Lin}(X)|$ is the ordered Bell number of $|X|$: $1, 1, 3, 13, 75, 541,\dots$
Eager linearisation is therefore expensive, which is exactly why \QED{}'s \emph{general}
algorithm (\S\ref{sec:qed-general}) defers permutation to the matching step instead.
\end{remark}

\subsubsection{Unification and completeness}
Terms are grouped by relational scope; within a group the summation disappears and only the
bodies are compared, because
\[
\usum{S}\ind{P_1}+\cdots+\ind{P_n} \;=\; \usum{S}\ind{P'_1}+\cdots+\ind{P'_m}
\quad\Longleftarrow\quad
\vec s \vdash_{\mathcal O} \ind{P_1}+\cdots+\ind{P_n} = \ind{P'_1}+\cdots+\ind{P'_m}
\]
and the right-hand side is an ordinary SMT query ($\ind{\cdot}$ becomes \texttt{ite},
$+$ becomes integer addition).

\begin{theorem}[\cite{qed} \S4.4]
On $\mathcal F_T$, with a complete oracle for $T$,
$\den{Q_1}=\den{Q_2}\Rightarrow \textsf{unify}(\textsf{lin}(\textsf{snf}(Q_1)),\textsf{lin}(\textsf{snf}(Q_2)))$.
\end{theorem}
\begin{proof}[Proof idea --- worth understanding, it is elegant]
Contrapositive. Suppose some group of terms with relational scope $S^*$ fails its body
check; choose such $S^*$ with the fewest variables. Build a database instance in which each
table $R_i$ is supported on exactly the finitely many points that the failing assignment
$\sigma$ touches, with multiplicity $n_i$ a free parameter. Because the term is in LNF, all
variables attached to a table are distinct and ordered, so a group with scope
$R_1^{k_1}\cdots R_m^{k_m}$ evaluates to the \emph{monomial}
$B_S[\delta,\sigma]\times n_1^{k_1}\cdots n_m^{k_m}$. Terms with larger scopes vanish (some
variable is forced outside the support). Both sides are thus polynomials in
$n_1,\dots,n_m$, and two polynomials over $\mathbb{N}$ are equal iff all coefficients match
--- but the $S^*$ coefficient differs by assumption. Hence some instantiation separates
$Q_1$ and $Q_2$.
\end{proof}

\keybox{The idea to steal}{%
LNF turns each term into a monomial whose exponents are the scope and whose coefficient is
the body. Bag equivalence then becomes \emph{polynomial identity}. That is why the paper
insists on a total order on summation variables: without it the exponents are not
well-defined.}

\subsection{Algorithm B: the general algorithm}
\label{sec:qed-general}

For everything outside $\mathcal F_T$ --- including our running example, which has
\texttt{GroupBy} --- \QED{} uses a weaker but far more general pipeline:
\textbf{normalise $\to$ stabilise $\to$ unify}. Here the oracle is allowed to answer
``unknown''.

\subsubsection{SPNF (\cite{qed} Alg.\ 4)}
Terms take the shape $\usum{\vec s}\ind{L}\times V$. The conversion is structural, and the
two interesting lines are
\[
\textsf{spnf}(\sq{V}) = \ind{\textsf{spnf}(V)\neq 0},
\qquad
\textsf{spnf}(\notu{V}) = \ind{\textsf{spnf}(V)=0},
\]
i.e.\ squash and negation are \emph{eliminated} in favour of predicates. Products merge
scopes and conjoin the conditions; sums concatenate.

\subsubsection{Stabilisation (\cite{qed} Alg.\ 5)}
\label{sec:stabilise}
This is \QED{}'s replacement for \udp{}'s Lemma~\ref{lem:15}, and the upgrade is that it is
\emph{semantic}. For each term $\usum{\vec s}\ind{L}\times V$ and each variable $s_i$, ask
whether $L$ implies $s_i = f(\text{other variables})$ for some expression $f$; if so,
substitute and delete the summation.

\begin{lstlisting}[style=pseudo]
function stable(U):
  U' <- 0
  for each term (Sum over s) [L] x V of U:
    s',L',V' <- s,L,V
    for each variable s_i in s:
      s'' <- s' minus s_i
      if exists f with  Gamma|Delta,s' |Phi |-_O  L' -> f(s'') = s_i :
         s' <- s'' ;  L' <- L'[f(s'')/s_i] ;  V' <- V'[f(s'')/s_i]
    U' <- U' + (Sum over s') [L'] x V'
  return U'
\end{lstlisting}

Two ways to find $f$ are given: synthesise it with SyGuS, or ask the SMT solver for the
\emph{congruence classes} of the term's subexpressions and read off an equal expression.
\QED{} implements the second for simplicity and reports the two are equally powerful in
practice.

\begin{example}[composing two projections]
$\usum{y,z}\ind{x=f(y)\wedge y=g(z)}R(z) \stackrel{?}{=} \usum{z}\ind{x=f(g(z))}R(z)$.
Congruence classes on the left over $\{x,y,z,f(y),g(z)\}$ are $\{x,f(y)\},\{y,g(z)\},\{z\}$,
exposing $y=g(z)$, so $y$ is eliminated. On the right the classes are $\{x,f(g(z))\},\{z\}$
and nothing is eliminated. The scopes now match.
\end{example}

\begin{remark}[a real, acknowledged gap]
The congruence-class method finds only expressions \emph{already present} in the term.
Given $\ind{x+3=10}$ it will not discover $x=7$, because $7$ never appears; SyGuS would.
The paper reports this never mattered in the evaluation.
\end{remark}

\subsubsection{Unification (\cite{qed} Alg.\ 6)}
\begin{lstlisting}[style=pseudo]
function unify(U1, U2):
  drop terms whose body is provably false (oracle: not (L and squash V))
  if #terms(U1) != #terms(U2): return Unknown       # <-- ceiling of THIS algorithm
  for each term T1 of U1:
     find a term T2 of U2 with termEq(T1,T2); remove it; else return Unknown
  return True

function termEq( (Sum s1)[L1] x V1 , (Sum s2)[L2] x V2 ):
  if len(s1) != len(s2): return False
  for each bijection sigma from s2 to s1:            # capped at 24 in the implementation
     L2,V2 <- L2[sigma], V2[sigma]
     L1',V1',L2',V2' <- ... with every HOp/QOp replaced by a FRESH variable
     C <- congruence facts obtained by recursively unifying those operators' arguments
          under the extra assumption  Phi and (L1' or L2')
     if  Phi and C  |-_O  [L1'] x V1' = [L2'] x V2' : return True
  return False
\end{lstlisting}

Three mechanisms are packed in here; take them one at a time.

\paragraph{(i) Scope permutation.} Summation is commutative, so before comparing bodies you
must try bijections between the two scopes. The justification is
$(\forall \vec s.\ f(\vec s)=g(\vec s)) \Rightarrow \usum{\vec s}f=\usum{\vec s}g$. Note
this is \emph{pointwise} comparison: sound, but it can only ever prove sums equal by proving
their summands equal. The implementation caps at 24 permutations; empirically only $\sim$2\%
of cases need more than the identity permutation.

\paragraph{(ii) Higher-order operators become fresh variables.}
$\qop(o,\vec v,\lam{x}{W})$ becomes a fresh table symbol; $\hop(o,\vec v,Q)$ a fresh scalar.
Congruence facts like $\forall x.R(x)=R'(x)$ are then \emph{recovered} by checking that the
operators match and recursively comparing their arguments.

\paragraph{(iii) The ambient assumption --- the genuinely new idea.}
\label{sec:ambient}
The recursive comparison of the operators' arguments is performed under the extra hypothesis
$L'_1\vee L'_2$: the conditions of the two enclosing terms. Soundness in one sentence:

\begin{lemma}[\cite{qed} Eq.\ 17]
To verify $\ind{L'_1}\times V'_1=\ind{L'_2}\times V'_2$ it suffices to check
$\Phi\wedge(L'_1\vee L'_2)\vdash C$ and $\Phi\wedge C\vdash \ind{L'_1}V'_1=\ind{L'_2}V'_2$.
\end{lemma}
\begin{proof}
If $\neg L'_1\wedge\neg L'_2$ then $\ind{L'_1}=\ind{L'_2}=0$ and both sides are $0$
regardless of $C$. So assuming $L'_1\vee L'_2$ while deriving $C$ loses nothing.
\end{proof}

\subsubsection{Encoding to first-order logic}
$\Nbar$-valued expressions cannot be handed to an SMT solver directly, but \QED{} never
needs their values --- only whether they are zero. So $U\neq0$ is pushed inward:
\[
U_1+U_2\neq0 \rightsquigarrow U_1\neq0\vee U_2\neq0,\qquad
U_1\times U_2\neq0 \rightsquigarrow U_1\neq0\wedge U_2\neq0,
\]
\[
\sq{U}\neq0\rightsquigarrow U\neq0,\qquad
\ind{P}\neq0\rightsquigarrow P,\qquad
\usum{s}U\neq0\rightsquigarrow \exists s.\,U\neq0.
\]
What survives is only fully-applied table symbols $R(a)\neq0$ and operator applications, all
of which are uninterpreted predicates. \QED{} additionally assumes every table has
\emph{finite} multiplicity at each row (type $S\to\mathbb{N}$, not $S\to\Nbar$) so that no
$\infty$ arises during term comparison.

\warnbox{This trick only works under a squash}{%
Turning $\usum{}$ into $\exists$ requires that the enclosing context care about
zero-vs-non-zero. When the top-level expression is \emph{not} squashed --- plain
\lstinline|SELECT x FROM r|, an unsquashed join --- multiplicities matter and \QED{} must
fall back on pointwise term matching, which is where its incompleteness bites
(\S\ref{sec:compare}).}

\begin{remark}[the two algorithms differ exactly here --- do not conflate them]
Alg.~6 demands a \emph{bijection} between the two sides' terms. Alg.~3, the
complete-fragment algorithm of \S\ref{sec:qed-complete}, does \textbf{not}: it groups terms
by relational scope and compares $\ind{P_1}+\cdots+\ind{P_n}$ with
$\ind{P'_1}+\cdots+\ind{P'_m}$ for $n\neq m$. So ``\QED{} cannot handle differing numbers of
terms'' is a statement about Alg.~6 only. A \lstinline|UNION ALL| split stays provable while
the query remains inside $\mathcal F_T$; adding a \texttt{GroupBy} or \texttt{Distinct}
pushes it into Alg.~6 and the bijection requirement returns.
\end{remark}

\subsection{The running example, proved}
\label{sec:qed-worked}

Translate $Q_1,Q_2$ of \S\ref{sec:running}, with the primary key already applied:
\begin{align*}
Q_1(x,r) &= \Bigl\lVert\usum{y,a}\ind{x{=}a}R(x,y)\sq{S(a)}\Bigr\rVert
 \times\Bigl[r=\Sumagg\bigl(\lam{y}{\usum{a}\ind{x{=}a}R(x,y)\sq{S(a)}}\bigr)\Bigr]\\
Q_2(x,r) &= \usum{a}\ind{x{=}a}\ \Bigl\lVert\usum{y}R(x,y)\Bigr\rVert
 \times\bigl[r=\Sumagg(\lam{y}{R(x,y)})\bigr]\times\sq{S(a)}
\end{align*}

\paragraph{Step 1 --- stabilise.} Both terms contain $\ind{x=a}$, which pins $a$ to $x$;
congruence classes put $x$ and $a$ together, so $\usum{a}$ is deleted:
\begin{align*}
Q_1(x,r) &= \Bigl\lVert\usum{y}R(x,y)\sq{S(x)}\bigl[r=\Sumagg(\lam{y}{R(x,y)\sq{S(x)}})\bigr]\Bigr\rVert\\
Q_2(x,r) &= \Bigl\lVert\usum{y}R(x,y)\sq{S(x)}\bigl[r=\Sumagg(\lam{y}{R(x,y)})\bigr]\Bigr\rVert
\end{align*}
One difference remains: the $\sq{S(x)}$ inside the aggregate.

\begin{remark}[why $Q_2$ may be written with an outer squash]
Every factor of $Q_2$ is already $0$-or-$1$: $\sq{\usum{y}R(x,y)}$, the bracket, and
$\sq{S(a)}$. So wrapping in $\sq{\cdot}$ changes nothing. But $\sq{S(a)}$ is squashed
\emph{only because of the primary key} --- drop the key and the factor is $S(a)$, which can
be $2$, and $Q_2$ is no longer squashable. That is precisely the counterexample of
\S\ref{sec:pkcounter}, where $Q_2$ returned $(1,30)$ \emph{twice}.
\end{remark}

\paragraph{Step 2 --- encode.} With $v_1,v_2$ fresh for the two aggregates:
\[
\phi_1(x,r)=\exists y.\,R(x,y)\neq0\wedge S(x)\neq0\wedge r=v_1,\qquad
\phi_2(x,r)=\exists y.\,R(x,y)\neq0\wedge S(x)\neq0\wedge r=v_2
\]
Goal $\forall x,r.\ \phi_1\leftrightarrow\phi_2$; identical but for $v_1$ vs $v_2$.

\paragraph{Step 3 --- recursive check with the ambient assumption.}
$v_1=v_2$ needs $\lam{y}{R(x,y)\sq{S(x)}}\equiv\lam{y}{R(x,y)}$, which is \textbf{false}
standalone ($x\notin S \Rightarrow$ left is empty). Under $\phi_1\vee\phi_2$, however, we
have $S(x)\neq0$, hence $\sq{S(x)}=1$, hence
\[
(\phi_1\vee\phi_2)\rightarrow\forall y.\ R(x,y)\times\sq{S(x)}=R(x,y)
\qquad\text{\emph{valid}.}
\]
So $v_1=v_2$, so $\phi_1\leftrightarrow\phi_2$, so $Q_1\equiv Q_2$. $\blacksquare$

\medskip
\noindent In English: \emph{the group-by only emits a row for an $x$ that survived the join,
and any such $x$ is in $S$; so inside the aggregate $\sq{S(x)}$ is just $1$ and vanishes.}

\subsection{Results and honest limits}

\begin{table}[h]\centering\small
\begin{tabular}{@{}llrrrrr@{}}
\toprule
suite & & \udp{} & EQUITAS & \spes{} & $\QED'$ & \QED{}\\
\midrule
Calcite (old, 232)   & proved & 34 & 67 & 95 & 49 & \textbf{147}\\
Calcite (v1.32, 444) & proved & -- & -- & 121 & 98 & \textbf{299}\\
CockroachDB (1287)   & proved & -- & -- & 325 & 656 & \textbf{979}\\
\bottomrule
\end{tabular}
\caption{\cite{qed} Table 1. $\QED'$ is the complete-fragment algorithm alone; \QED{} is the
full pipeline. Every case \spes{} proves, \QED{} also proves. Average runtime
0.88\,s (Calcite v1.32), of which $\sim$0.83\,s is SMT time inside stabilisation.}
\end{table}

\begin{table}[h]\centering\small
\begin{tabular}{@{}lrr@{}}
\toprule
failure category & Calcite & CockroachDB\\
\midrule
aggregates        & \textbf{96} & 27\\
custom operators  & 25 & \textbf{152}\\
list semantics (\lstinline|ORDER BY|/\lstinline|LIMIT|) & 18 & 103\\
typing (implicit casts) & 6 & 26\\
\bottomrule
\end{tabular}
\caption{\cite{qed} Table 3. Aggregates dominate on Calcite --- the direct cost of leaving
$\alpha$ uninterpreted. \sqlsolver{} attacks exactly this column.}
\end{table}
